\chapter{Source-Event Reporting Prototype}
\label{ch:ground-compiler-source-event-reporting-prototype}

\autoref{ch:ground-compiler-minimal-executable-prototype} defines the minimal executable prototype $P_0$, which encodes event flows, decodes channel streams, checks legality, rejects illegal inputs, and runs round-trip tests. This chapter defines the next prototype level:
$$
  \boxed{P_1=P_0+\text{decoded source-event reports}.}
$$
The $P_1$ layer is a source-event reporting prototype. It does not recognize packages, name certificates, derivation certificates, theorems, chapters, or accepted objects. It decodes a legal channel stream into an event flow and emits a structured human-readable report.

The guiding rule is:
$$
  \boxed{\text{a report is an output view, not an input, not a recognizer, and not a proof}.}
$$

\section{Source-event reports}
\label{sec:source-report-reports}

\begin{definition}[Source-event reporting prototype]
\label{def:source-report-prototype}
A \emph{source-event reporting prototype} is an executable prototype $P_1$ satisfying:
\begin{enumerate}
\item $P_1$ includes the encode, decode, check, and round-trip functions of $P_0$;
\item if the input is a legal channel stream $c$, then $P_1$ decodes
$$
  S=\mathsf{Decode}(c);
$$
\item $P_1$ emits a source-event report $\mathsf{Report}_{\mathrm{src}}(c,S)$;
\item the report describes source events, channel segments, round-trip status, warnings, and cannot-claim annotations;
\item the report does not recognize packages, name certificates, theorems, chapters, or accepted objects.
\end{enumerate}
Write
$$
  \mathsf{SourceReportPrototype}(P_1).
$$
\end{definition}

\begin{principle}[Report is output, not input]
\label{prin:source-report-output-not-input}
Reports generated by $P_1$ are output views. They cannot be used as formal compiler input structures. In particular, there is no formal inference step from report to package, name certificate, theorem, or accepted object. The allowed direction is:
$$
  \text{channel stream}
  \longrightarrow
  \mathsf{Decode}
  \longrightarrow
  \text{event flow}
  \longrightarrow
  \text{source-event report}.
$$
\end{principle}

\begin{definition}[Source-event report]
\label{def:source-report-source-event-report}
Given a legal channel stream
$$
  c\in\mathsf{LegalZStream}
$$
and decoded event flow
$$
  S=\mathsf{Decode}(c)=(w_0,\ldots,w_{n-1}),
$$
a \emph{source-event report} is an output view containing at least the input channel stream, legality status, decoded event count, event index, raw event, channel segment for each event, body segment for each event, terminator segment for each event, round-trip check, source/channel separation warning, and cannot-claim annotations.

Write
$$
  \mathsf{SrcReport}(c,S).
$$
\end{definition}

\begin{definition}[Event segment]
\label{def:source-report-event-segment}
\leandef{BEDC.GroundCompiler.SourceReport.EventSegment}
Let $S=(w_0,\ldots,w_{n-1})$ and $c=\mathsf{Compile}(S)$. The \emph{event segment} of the $i$-th event is
$$
  \mathsf{Seg}_i(c,S):=\mathsf{Ev}(w_i).
$$
This is a channel-level substring corresponding to a complete event encoding, not an arbitrary channel substring.
\end{definition}

\begin{definition}[Body and terminator segment]
\label{def:source-report-body-terminator-segment}
\leandef{BEDC.GroundCompiler.SourceReport.BodyTerminatorSegments}
For an event $w_i$, its event segment is
$$
  \mathsf{Seg}_i=\zeta(w_i)11.
$$
Define
$$
  \mathsf{Body}_i:=\zeta(w_i),
  \qquad
  \mathsf{Term}_i:=11.
$$
Then
$$
  \mathsf{Seg}_i=\mathsf{Body}_i\mathsf{Term}_i.
$$
\end{definition}

\begin{proposition}[Event segments partition the channel stream]
\label{prop:source-report-segments-partition}
\leanchecked{BEDC.GroundCompiler.SourceReport.event\_segments\_partition}
If
$$
  c=\mathsf{Compile}(w_0,\ldots,w_{n-1}),
$$
then
$$
  c=\mathsf{Seg}_0\mathsf{Seg}_1\cdots\mathsf{Seg}_{n-1}.
$$
\end{proposition}

\begin{proof}
By the flow compiler definition,
$$
  \mathsf{Compile}(w_0,\ldots,w_{n-1})
  =
  \mathsf{Ev}(w_0)\cdots\mathsf{Ev}(w_{n-1}).
$$
By definition, $\mathsf{Seg}_i=\mathsf{Ev}(w_i)$. Substitution gives the displayed partition.
\end{proof}

\begin{corollary}[Reports may display event segmentation]
\label{cor:source-report-may-display-segments}
Since event segments partition the channel stream, $P_1$ may display segment $0$ through segment $n-1$ without introducing an external tree or AST. The segmentation is induced by channel decoding.
\end{corollary}

\begin{definition}[Event index]
\label{def:source-report-event-index}
For a decoded event flow $S=(w_0,\ldots,w_{n-1})$, the \emph{event index} of $w_i$ is $i$. In a report, event index is implementation-level display. It is not a BEDC source event and not part of the formal event flow.
\end{definition}

\begin{remark}[Event index is not source structure]
\label{rem:source-report-index-not-source-structure}
The index $i$ helps a human inspect the decoded flow. It is not a raw event, theorem label, package label, or object identity. A report may print $w_i=011$, but the formal source event is just $011$.
\end{remark}

\section{Warnings and cannot-claims}
\label{sec:source-report-warnings-cannot-claims}

\begin{definition}[Literal source warning]
\label{def:source-report-literal-warning}
A \emph{literal source warning} is a report entry generated from raw source events without claiming motif recognition. Examples include: an event $011$ followed by $100$, an event containing adjacent source $11$, an event resembling a raw completion witness, or an event resembling a normal-address candidate.

Every literal warning must state that it is candidate-only and requires recognizer flow for interpretation.
\end{definition}

\begin{principle}[Literal warning is not motif recognition]
\label{prin:source-report-warning-not-recognition}
A literal source warning is not a motif recognition judgment. Thus a warning that $011$ is followed by $100$ does not imply a Zeckendorf-carry relation and does not imply a classifier equivalence. It only states that the decoded event flow contains the literal adjacent events $011\mid 100$.
\end{principle}

\begin{definition}[Cannot-claim annotation]
\label{def:source-report-cannot-claim-annotation}
A \emph{cannot-claim annotation} in a source-event report records what the report does not establish. Required examples include:
\begin{itemize}
\item decoded event flow is not theoremhood;
\item decoded event flow is not a name certificate;
\item decoded event flow is not a derivation certificate;
\item decoded event flow is not an acceptance gate;
\item literal $011\mid 100$ is not carry without recognizer;
\item channel terminator $11$ is not source event $11$;
\item report is not proof;
\item report is not package recognition.
\end{itemize}
\end{definition}

\begin{principle}[Nontrivial reports include cannot-claims]
\label{prin:source-report-nontrivial-cannot-claims}
If a report displays events commonly used as mathematical witnesses, such as $01$, $011$, $100$, $001$, $0011$, or $0100$, then it must include cannot-claim annotations preventing overinterpretation.
\end{principle}

\begin{definition}[Source report for illegal input]
\label{def:source-report-illegal-input}
If channel input $c$ is illegal, $P_1$ does not produce a source-event report. It produces a rejection report containing the input, failure position, rejection reason, optionally a partial decoded prefix if policy allows it, and a cannot-claim annotation saying that an illegal channel stream has no decoded event flow.
\end{definition}

\begin{proposition}[Illegal channel stream has no source-event report]
\label{prop:source-report-illegal-no-report}
If $c\notin\mathsf{LegalZStream}$, then $P_1$ cannot output $\mathsf{SrcReport}(c,S)$ for any $S$.
\end{proposition}

\begin{proof}
A source-event report requires $S=\mathsf{Decode}(c)$. But $\mathsf{Decode}$ is defined on legal channel streams. If $c$ is not legal, no decoded event flow $S$ is available.
\end{proof}

\begin{definition}[Round-trip report]
\label{def:source-report-roundtrip-report}
A \emph{round-trip report} for input event flow $S$ records $S$, $\mathsf{Compile}(S)$, $\mathsf{Decode}(\mathsf{Compile}(S))$, and the comparison result. It succeeds when
$$
  \mathsf{Decode}(\mathsf{Compile}(S))=S.
$$
For input channel stream $c$, it records $c$, $\mathsf{Decode}(c)$, $\mathsf{Compile}(\mathsf{Decode}(c))$, and the comparison result. It succeeds when
$$
  \mathsf{Compile}(\mathsf{Decode}(c))=c.
$$
\end{definition}

\begin{proposition}[Round-trip report is determined by channel bijection]
\label{prop:source-report-roundtrip-determined}
If $S\in\mathsf{EventFlow}$, then its round-trip report succeeds. If $c\in\mathsf{LegalZStream}$, then its round-trip report succeeds.
\end{proposition}

\begin{proof}
This is exactly the channel bijection theorem:
$$
  \mathsf{Decode}(\mathsf{Compile}(S))=S,
  \qquad
  \mathsf{Compile}(\mathsf{Decode}(c))=c.
$$
\end{proof}

\section{Report format and examples}
\label{sec:source-report-format-examples}

\begin{definition}[Minimal source-event report format]
\label{def:source-report-minimal-format}
A minimal report contains input channel stream, legal status, decoded event-flow display when legal, event index, raw event, body code, terminator code, full event code, round-trip status, literal source warnings, and cannot-claim annotations.
\end{definition}

\begin{example}[Report for completion skeleton]
\label{ex:source-report-completion-skeleton}
For input
$$
  011001100011010110101011100011,
$$
the decoded event flow is
$$
  0\mid 00\mid 000\mid 01\mid 011\mid 100.
$$
The event table records:
$$
\begin{array}{c|c|c|c|c}
\text{index} & \text{raw} & \text{body} & \text{terminator} & \text{code}\\
\hline
0 & 0   & 0     & 11 & 011\\
1 & 00  & 00    & 11 & 0011\\
2 & 000 & 000   & 11 & 00011\\
3 & 01  & 010   & 11 & 01011\\
4 & 011 & 01010 & 11 & 0101011\\
5 & 100 & 1000  & 11 & 100011.
\end{array}
$$
Warnings include that the literal pair $011\mid 100$ appears, but this is not carry recognition and requires source-level carry recognizer flow. Cannot-claims include that the decoded flow does not license $\FullAxisUp$, Completion, or dimension-lift structure, that $011$ is not equal to $100$, that channel terminator $11$ is not source event $11$, and that the report is not proof.
\end{example}

\begin{example}[Report for Add/fold skeleton]
\label{ex:source-report-add-fold-skeleton}
For input
$$
  0110011001011001010110100011,
$$
the decoded event flow is
$$
  0\mid 00\mid 001\mid 0011\mid 0100.
$$
The event table records raw events $0$, $00$, $001$, $0011$, and $0100$ with codes $011$, $0011$, $001011$, $00101011$, and $0100011$. The report may warn that the literal pair $0011\mid 0100$ appears, but this is not fold recognition without recognizer flow. Cannot-claims include that $001$ is not $\AddUp$, $0011$ is not $\MulUp$, $0100$ is not an accepted fold object, and the report supplies no name-certificate or derivation-certificate flow.
\end{example}

\begin{example}[Report for raw source event $11$]
\label{ex:source-report-raw-source-11}
For input
$$
  101011,
$$
the decoded event flow is the single source event $11$. The event table records raw $11$, body $1010$, terminator $11$, and code $101011$. Warnings state that source raw event $11$ is represented in channel as $101011$, that the final $11$ is the channel terminator, and that the source $11$ is encoded as $10\,10$. Cannot-claims state that channel terminator is not source $11$, source $11$ is not automatically a seal, and source $11$ is not automatically completion.
\end{example}

\section{Soundness and display policy}
\label{sec:source-report-soundness-display}

\begin{definition}[Sound source report]
\label{def:source-report-sound}
A source-event report $\mathsf{SrcReport}(c,S)$ is \emph{sound} when:
\begin{enumerate}
\item $c\in\mathsf{LegalZStream}$;
\item $\mathsf{Decode}(c)=S$;
\item every event-table row corresponds to an actual event $w_i$ of $S$;
\item every displayed event code equals $\mathsf{Ev}(w_i)$;
\item every warning is marked candidate-only unless backed by recognizer flow;
\item every cannot-claim annotation is consistent with the limited capability of $P_1$.
\end{enumerate}
\end{definition}

\begin{theorem}[Sound report theorem]
\label{thm:source-report-sound-report}
If $P_1$ produces a sound source-event report $\mathsf{SrcReport}(c,S)$, then
$$
  \mathsf{Decodes}(c,S).
$$
\end{theorem}

\begin{proof}
By soundness of the report, $c\in\mathsf{LegalZStream}$ and $\mathsf{Decode}(c)=S$. Hence $\mathsf{Decodes}(c,S)$.
\end{proof}

\begin{corollary}[Source report preserves raw event flow]
\label{cor:source-report-preserves-raw-flow}
A sound source-event report records exactly the decoded source event flow. It does not add or remove source events.
\end{corollary}

\begin{theorem}[Source report does not recognize structures]
\label{thm:source-report-does-not-recognize-structures}
A sound $P_1$ source-event report does not establish package, name-certificate, derivation-certificate, theorem, chapter, motif, or accepted-object recognition.
\end{theorem}

\begin{proof}
The $P_1$ layer decodes channel streams and reports source events. It has no generated recognizer flow for package, certificate, theorem, chapter, motif, or accepted-object roles. Therefore it cannot establish those recognitions.
\end{proof}

\begin{corollary}[Warnings are weaker than recognitions]
\label{cor:source-report-warnings-weaker}
A warning such as ``literal pair $011\mid 100$ appears'' is not carry motif recognition. It only reports a literal pattern in decoded source events.
\end{corollary}

\begin{theorem}[Cannot-claims prevent overinterpretation]
\label{thm:source-report-cannot-claims-prevent-overinterpretation}
If a $P_1$ source report contains raw events commonly used as higher-layer witnesses, then cannot-claim annotations are required to prevent the report from being read as proof or certificate.
\end{theorem}

\begin{proof}
At $P_1$, no higher recognizers are present. Any interpretation of raw events as completion, carry, theorem, package, or accepted object would exceed $P_1$ capability. Cannot-claim annotations record this limitation explicitly. Without them, the report could be mistaken for a stronger claim.
\end{proof}

\begin{definition}[Display policy]
\label{def:source-report-display-policy}
A \emph{display policy} is an implementation-level convention determining how event flows and reports are printed. It may specify the raw event separator, empty-event marker, event-index format, indentation, warning format, and cannot-claim format. Display policy is output formatting, not formal input.
\end{definition}

\begin{principle}[Display policy is not semantics]
\label{prin:source-report-display-policy-not-semantics}
Changing display policy must not change $\mathsf{Decode}(c)$ or $\mathsf{Compile}(S)$. Display policy may change how reports look, but not what event flow they report.
\end{principle}

\begin{proposition}[Display policy does not affect round trip]
\label{prop:source-report-display-policy-roundtrip}
If two reports differ only by display policy, they represent the same decoded event flow and the same channel stream.
\end{proposition}

\begin{proof}
Display policy affects only output formatting. It does not alter $\mathsf{Compile}$ or $\mathsf{Decode}$.
\end{proof}

\begin{definition}[Canonical display policy]
\label{def:source-report-canonical-display-policy}
A canonical display policy for $P_1$ may use separator ``$\mid$'', empty event marker $\epsilon$, event index form $i$, explicit body, terminator, and code fields, warnings after the event table, and cannot-claims last. This policy is recommended for reproducibility but remains output policy, not source semantics.
\end{definition}

\section{No-host-leak and adequacy}
\label{sec:source-report-no-host-leak-adequacy}

\begin{definition}[Host leak in report]
\label{def:source-report-host-leak}
A source report contains a \emph{host leak} if it presents host-level structures as formal BEDC structures. Examples include a Python list interpreted as an event-flow object, JSON object interpreted as a package, YAML object interpreted as a chapter, string label interpreted as theorem identity, or host enum interpreted as status flow.
\end{definition}

\begin{theorem}[Host-leaking report is inadmissible]
\label{thm:source-report-host-leak-inadmissible}
If a $P_1$ report contains a host leak, then it is not an admissible source-event report.
\end{theorem}

\begin{proof}
A report is an output view of decoded event flow. Presenting host structures as formal BEDC structures introduces hidden input or misrepresents implementation artifacts as source-level objects. This violates the no-host-leak condition.
\end{proof}

\begin{corollary}[JSON and YAML reports are views only]
\label{cor:source-report-json-yaml-view-only}
JSON or YAML output is allowed only when clearly marked as a human-readable report. It cannot be re-imported as formal compiler input.
\end{corollary}

\begin{definition}[$P_1$ audit checklist]
\label{def:source-report-p1-audit}
A $P_1$ prototype must satisfy:
\begin{enumerate}
\item all $P_0$ audit requirements;
\item legal channel streams decode to event flows;
\item reports display exact event segmentation;
\item reports display body and terminator separately;
\item reports flag source/channel distinctions;
\item reports include cannot-claim annotations for high-risk events;
\item reports do not claim package, theorem, certificate, or accepted-object recognition;
\item reports do not use host structures as formal objects;
\item report round trips agree with $\mathsf{Compile}$ and $\mathsf{Decode}$;
\item illegal input generates rejection report, not source report.
\end{enumerate}
\end{definition}

\begin{theorem}[$P_1$ adequacy theorem]
\label{thm:source-report-p1-adequacy}
If a prototype satisfies the $P_1$ audit checklist and its reports are sound, then it is adequate as a source-event reporting prototype.
\end{theorem}

\begin{proof}
The checklist ensures that the prototype implements all $P_0$ functions and faithfully reports decoded source events. Sound reports guarantee that displayed event flows match $\mathsf{Decode}(c)$. Cannot-claim annotations prevent overinterpretation, and no-host-leak restrictions prevent reports from becoming hidden input. Therefore the prototype is adequate at $P_1$.
\end{proof}

\begin{corollary}[$P_1$ adequacy does not imply higher adequacy]
\label{cor:source-report-p1-not-higher}
A $P_1$ prototype is not a source-level normalizer, package recognizer, name-certificate recognizer, theorem recognizer, motif analyzer, or self-hosting compiler.
\end{corollary}

\begin{theorem}[Source-event reporting is conservative]
\label{thm:source-report-conservative}
Adding $P_1$ source-event reports adds no raw marks, history constructors, sameness rules, continuation rules, derived-interface formation rules, theorem rules, or accepted-object export rules.
\end{theorem}

\begin{proof}
$P_1$ only decodes channel streams into event flows and prints reports. It introduces no new source events, asserts no $\msame$, $\hsame$, or $\psame$, modifies neither $\ExtR$ nor $\ContR$, and recognizes or licenses no higher structure. Hence it is conservative over the finite kernel.
\end{proof}

\begin{corollary}[$P_1$ remains an address/report layer]
\label{cor:source-report-p1-address-report-layer}
$P_1$ remains an address and reporting layer. It is not a mathematical object generator.
\end{corollary}

\section{Full example report}
\label{sec:source-report-full-example}

\begin{example}[Full $P_1$ report for completion skeleton]
\label{ex:source-report-full-completion}
For input
$$
  011001100011010110101011100011,
$$
the report has legal status true and decoded event flow
$$
  0\mid 00\mid 000\mid 01\mid 011\mid 100.
$$
Its event table is
$$
\begin{array}{c|c|c|c|c}
\text{index} & \text{raw} & \text{body} & \text{terminator} & \text{code}\\
\hline
0 & 0   & 0     & 11 & 011\\
1 & 00  & 00    & 11 & 0011\\
2 & 000 & 000   & 11 & 00011\\
3 & 01  & 010   & 11 & 01011\\
4 & 011 & 01010 & 11 & 0101011\\
5 & 100 & 1000  & 11 & 100011.
\end{array}
$$
The round trip records
$$
  \mathsf{Decode}(c)=0\mid 00\mid 000\mid 01\mid 011\mid 100
$$
and
$$
  \mathsf{Compile}(\mathsf{Decode}(c))=c.
$$
Warnings record that the literal event pair $011\mid 100$ appears, with candidate interpretation only: possible source-level carry ledger, requiring generated carry recognizer flow, carry ledger, and classifier exactness. Cannot-claims record that this report does not prove $011=100$, does not prove $\hsame(011,100)$, does not license $\FullAxisUp$, Completion, or dimension-lift structure, and is not a name-certificate flow, theorem flow, or acceptance-gate flow.
\end{example}

\begin{remark}[Summary]
\label{rem:source-report-summary}
This chapter defines $P_1$, the source-event reporting prototype. It extends $P_0$ with decoded event tables, event segmentation reports, body/terminator separation, round-trip reports, literal source warnings, cannot-claim annotations, display policy, and host-leak restrictions. It proves that illegal channel streams have no source-event report, round-trip reports are determined by the channel bijection, sound reports imply decoding, reports preserve raw event flows, reports do not recognize higher structures, warnings are weaker than recognitions, cannot-claims prevent overinterpretation, display policy is not semantics, host-leaking reports are inadmissible, $P_1$ adequacy follows from audit and report soundness, and source-event reporting is conservative.
\end{remark}
